\documentclass[11pt]{article}
\usepackage[margin=1.1in]{geometry}
\usepackage{amsmath,amssymb,amsthm,mathtools,booktabs,array,tabularx,longtable}
\usepackage[dvipsnames]{xcolor}
\usepackage[colorlinks=true,linkcolor=MidnightBlue,citecolor=MidnightBlue,urlcolor=MidnightBlue]{hyperref}
\usepackage{fancyhdr}
\pagestyle{fancy}\fancyhf{}\renewcommand{\headrulewidth}{0pt}
\fancyfoot[R]{\thepage}
\fancyfoot[L]{\scriptsize\color{gray}Generated by the Transport Lab from approved parts of its store; not a deposited paper. Papers of record: \href{https://zenodo.org/communities/transport-papers}{Zenodo, Transport Papers}}
\fancypagestyle{plain}{\fancyhf{}\fancyfoot[R]{\thepage}\fancyfoot[L]{\scriptsize\color{gray}Generated by the Transport Lab from approved parts of its store; not a deposited paper. Papers of record: \href{https://zenodo.org/communities/transport-papers}{Zenodo, Transport Papers}}\renewcommand{\headrulewidth}{0pt}}
\newtheorem{theorem}{Theorem}[section]
\newtheorem{proposition}[theorem]{Proposition}
\newtheorem{lemma}[theorem]{Lemma}
\newtheorem{corollary}[theorem]{Corollary}
\newtheorem{conjecture}[theorem]{Conjecture}
\newtheorem{refutation}[theorem]{Refutation}
\newtheorem{prediction}[theorem]{Prediction}
\newtheorem{principle}[theorem]{Principle}
\newtheorem{observation}[theorem]{Observation}
\newtheorem{axiom}[theorem]{Axiom}
\theoremstyle{definition}
\newtheorem{definition}[theorem]{Definition}
\newtheorem{fact}[theorem]{Fact}
\newtheorem{assembled}[theorem]{Assembly}
\newtheorem{postulate}[theorem]{Postulate}
\newtheorem{premise}[theorem]{Premise}
\newtheorem{hypothesis}[theorem]{Hypothesis}
\newtheorem{remark}[theorem]{Remark}
\newtheorem{example}[theorem]{Example}
\theoremstyle{plain}
\setlength{\marginparwidth}{0.8in}
\newcommand{\changedmark}[1]{\marginpar{\raggedright\scriptsize\color{BrickRed}#1}}
\providecommand{\Aa}{}\renewcommand{\Aa}{\mathcal{A}}
\providecommand{\Ad}{}\renewcommand{\Ad}{\widetilde{\mathrm{Ad}}}
\providecommand{\Aev}{}\renewcommand{\Aev}{\mathcal{A}^{+}}
\providecommand{\Aod}{}\renewcommand{\Aod}{\mathcal{A}^{-}}
\providecommand{\B}{}\renewcommand{\B}{\mathrm{B}}
\providecommand{\C}{}\renewcommand{\C}{\mathbb{C}}
\providecommand{\CAR}{}\renewcommand{\CAR}{\mathrm{CAR}}
\providecommand{\CS}{}\renewcommand{\CS}{{\mathbb{S}^{+}}}
\providecommand{\CSm}{}\renewcommand{\CSm}{{\mathbb{S}^{-}}}
\providecommand{\Cl}{}\renewcommand{\Cl}{\mathrm{Cl}}
\providecommand{\Con}{}\renewcommand{\Con}{\mathcal{C}}
\providecommand{\Dl}{}\renewcommand{\Dl}{\Delta}
\providecommand{\EN}{}\renewcommand{\EN}{E_{N}}
\providecommand{\End}{}\renewcommand{\End}{\mathrm{End}}
\providecommand{\F}{}\renewcommand{\F}{\mathbb{F}}
\providecommand{\G}{}\renewcommand{\G}{\Gamma}
\providecommand{\HH}{}\renewcommand{\HH}{\mathbb{H}}
\providecommand{\IS}{}\renewcommand{\IS}{I\sigma_3}
\providecommand{\J}{}\renewcommand{\J}{\mathbf{J}}
\providecommand{\LH}{}\renewcommand{\LH}{L_{H}}
\providecommand{\Mat}{}\renewcommand{\Mat}{\mathrm{M}}
\providecommand{\Mo}{}\renewcommand{\Mo}{M^{\circ}}
\providecommand{\Mp}{}\renewcommand{\Mp}{M^{+}}
\providecommand{\Msun}{}\renewcommand{\Msun}{M_{\odot}}
\providecommand{\Mt}{}\renewcommand{\Mt}{\widetilde{M}}
\providecommand{\Ncone}{}\renewcommand{\Ncone}{\mathcal{N}}
\providecommand{\Neff}{}\renewcommand{\Neff}{N_{\mathrm{eff}}}
\providecommand{\PCS}{}\renewcommand{\PCS}{\mathbb{P}{\mathbb{S}^{+}}}
\providecommand{\PT}{}\renewcommand{\PT}{\mathbb{PT}}
\providecommand{\Pin}{}\renewcommand{\Pin}{\mathrm{Pin}}
\providecommand{\Q}{}\renewcommand{\Q}{\mathbb{Q}}
\providecommand{\R}{}\renewcommand{\R}{\mathbb{R}}
\providecommand{\Ric}{}\renewcommand{\Ric}{\mathrm{Ric}}
\providecommand{\SBH}{}\renewcommand{\SBH}{S_{\mathrm{BH}}}
\providecommand{\Sc}{}\renewcommand{\Sc}{\mathbb{S}}
\providecommand{\Sch}{}\renewcommand{\Sch}{\mathrm{P}}
\providecommand{\Sig}{}\renewcommand{\Sig}{\Sigma}
\providecommand{\Som}{}\renewcommand{\Som}{S_\omega}
\providecommand{\Sp}{}\renewcommand{\Sp}{\mathbb{S}}
\providecommand{\Sph}{}\renewcommand{\Sph}{\mathrm{S}}
\providecommand{\Spi}{}\renewcommand{\Spi}{S_\pi}
\providecommand{\Spin}{}\renewcommand{\Spin}{\mathrm{Spin}}
\providecommand{\Sth}{}\renewcommand{\Sth}{S_{\mathrm{th}}}
\providecommand{\T}{}\renewcommand{\T}{\mathbb{T}}
\providecommand{\TH}{}\renewcommand{\TH}{T_{\mathrm{H}}}
\providecommand{\Tr}{}\renewcommand{\Tr}{\operatorname{Tr}}
\providecommand{\UHF}{}\renewcommand{\UHF}{\mathrm{UHF}}
\providecommand{\Z}{}\renewcommand{\Z}{\mathbb{Z}}
\providecommand{\Zc}{}\renewcommand{\Zc}{\mathrm{Z}}
\providecommand{\cmg}{}\renewcommand{\cmg}{\mathrm{cm^{2}\,g^{-1}}}
\providecommand{\dS}{}\renewcommand{\dS}{\mathrm{dS}}
\providecommand{\dsc}{}\renewcommand{\dsc}{\mathrm{dsc}}
\providecommand{\dtPG}{}\renewcommand{\dtPG}{dt}
\providecommand{\dtS}{}\renewcommand{\dtS}{dt_{\mathrm{S}}}
\providecommand{\ellv}{}\renewcommand{\ellv}{\ell}
\providecommand{\er}{}\renewcommand{\er}{e_r}
\providecommand{\et}{}\renewcommand{\et}{e_t}
\providecommand{\ev}{}\renewcommand{\ev}{\mathrm{ev}}
\providecommand{\gl}{}\renewcommand{\gl}{\mathfrak{gl}}
\providecommand{\hatV}{}\renewcommand{\hatV}{\widehat{V}}
\providecommand{\hb}{}\renewcommand{\hb}{\bar{h}}
\providecommand{\hf}{}\renewcommand{\hf}{\hat{f}}
\providecommand{\id}{}\renewcommand{\id}{\mathrm{id}}
\providecommand{\kpc}{}\renewcommand{\kpc}{\mathrm{kpc}}
\providecommand{\lP}{}\renewcommand{\lP}{\ell_{P}}
\providecommand{\mD}{}\renewcommand{\mD}{\mathbf{m}_{D}}
\providecommand{\mb}{}\renewcommand{\mb}{\bar m}
\providecommand{\mnu}{}\renewcommand{\mnu}{\Sigma m_{\nu}}
\providecommand{\nInf}{}\renewcommand{\nInf}{n_\infty}
\providecommand{\nO}{}\renewcommand{\nO}{n_0}
\providecommand{\ninf}{}\renewcommand{\ninf}{n_{\infty}}
\providecommand{\nz}{}\renewcommand{\nz}{n_{0}}
\providecommand{\proj}{}\renewcommand{\proj}{\operatorname{proj}}
\providecommand{\pv}{}\renewcommand{\pv}{\pi}
\providecommand{\rad}{}\renewcommand{\rad}{\mathrm{rad}}
\providecommand{\sgn}{}\renewcommand{\sgn}{\operatorname{sgn}}
\providecommand{\shA}{}\renewcommand{\shA}{\mathcal{A}}
\providecommand{\shC}{}\renewcommand{\shC}{\mathcal{C}}
\providecommand{\shE}{}\renewcommand{\shE}{\mathcal{E}}
\providecommand{\shJ}{}\renewcommand{\shJ}{\mathcal{J}}
\providecommand{\shK}{}\renewcommand{\shK}{\mathcal{K}}
\providecommand{\shL}{}\renewcommand{\shL}{L}
\providecommand{\shW}{}\renewcommand{\shW}{\mathcal{W}}
\providecommand{\sll}{}\renewcommand{\sll}{\mathfrak{sl}}
\providecommand{\so}{}\renewcommand{\so}{\mathfrak{so}}
\providecommand{\su}{}\renewcommand{\su}{\mathfrak{su}}
\providecommand{\supp}{}\renewcommand{\supp}{\operatorname{supp}}
\providecommand{\tr}{}\renewcommand{\tr}{\operatorname{tr}}
\providecommand{\uu}{}\renewcommand{\uu}{\mathfrak{u}}
\providecommand{\br}{}\renewcommand{\br}[1]{\langle #1\rangle}
\providecommand{\bra}{}\renewcommand{\bra}[1]{\langle #1|}
\providecommand{\file}{}\renewcommand{\file}[1]{\texttt{#1}}
\providecommand{\hp}{}\renewcommand{\hp}[2]{( #1 \,|\, #2 )}
\providecommand{\ip}{}\renewcommand{\ip}[2]{\langle #1, #2\rangle}
\providecommand{\ket}{}\renewcommand{\ket}[1]{|#1\rangle}
\providecommand{\rel}{}\renewcommand{\rel}[2]{D\bigl(#1 \,\|\, #2\bigr)}
\providecommand{\rev}{}\renewcommand{\rev}[1]{\widetilde{#1}}
\providecommand{\slashed}{}\renewcommand{\slashed}[1]{{#1}\mkern-10mu/}
\title{\bfseries Three Routes from the Localisation: The Metric-Level Sheaf, the Boundary No-Go, and the Hopping Exclusion: A companion note to Transport Papers T5 (Steps 4a, 5 and 6b) and T4 (Section 8)}
\author{Leon Butler\\ \normalsize\href{mailto:lbutler.math@gmail.com}{\texttt{lbutler.math@gmail.com}}\\ \normalsize ORCID: \href{https://orcid.org/0009-0006-4449-5617}{0009-0006-4449-5617}}
\date{Generated by the Transport Lab, 27 October 2026}
\begin{document}
\maketitle
\begin{abstract}
\normalsize
Three open items of the Transport Papers are addressed at the point where they were priced. First, the pointwise reading of the metric level: with collar data in place, the assignment sending an open piece of the singular boundary to the families of fibrewise parity-even elements is a genuine sheaf, its existence as an untwisted sheaf equivalent to the standing hypothesis $w_1(\shL) = 0$; its stalks are the pointwise metric-level algebras, its finite truncations are exactly the range of the multi-mode descent, and the rival reading, the algebra generated by the pointwise fixed algebras, has the wrong stalks and is the section object of no sheaf. The antecedent of T5's Proposition 9.6 thereby reduces to the locality of the metric level. Second, the boundary no-go: completions built from different collar transversals are canonically isomorphic over the degenerate base, by an explicit unipotent intertwiner whose closure uses the radical gauge freedom; the base unipotents extend to automorphisms of any fixed completion and act on the candidate classical values by translations, which have no fixed point. No equivariant classical partner value exists at the boundary itself. Third, the hopping exclusion: every generator the tame machinery supplies is function-linear, function-linear endomorphisms preserve supports, and support-preserving families second-quantise to sums of single-point bilinears; occupation hopping between points is therefore impossible within the supplied class, the true residual being a genuinely differential dynamics. Every computational identity is machine-verified in exact arithmetic.
\end{abstract}
\medskip
\noindent\textbf{MSC 2020.} Primary 15A66, 53C50; Secondary 46L60, 18F20.\\
\textbf{Keywords.} Clifford algebra; degenerate metric; sheaf; conditional expectation; superselection; no-go theorem.
\newpage
\noindent \textbf{Authorship and disclosure. }The founding idea of the series is the author's: that division by zero in a Clifford algebra should be read not as a prohibition but as a transport into a higher-dimensional algebra, and that the transported division need not be accompanied by an inverse operation. The direction of investigation and the critical review of the results are likewise the author's. The mathematical development and the drafting of the text were carried out in extended collaboration with Claude Fable 5 (Anthropic), an AI system.

\noindent\fbox{\parbox{\dimexpr\linewidth-2\fboxsep-2\fboxrule}{\textbf{This is not a deposited paper.} It was generated by the Transport Lab on 27 October 2026 and assembled from approved parts of its store. It follows results across several papers of the series; each is stated with its proof and credited to the paper that first proved it. The official versions of those papers are on Zenodo, in the Transport Papers community, \url{https://zenodo.org/communities/transport-papers}.}}
\newpage

\section{Setting, and what is imported}

\phantomsection\label{c:note.three-routes.setting-setting-what-imported}
The habitat is T4's Schwarzschild composite \cite{ButlerT4}, in the form T5 reproduces it verbatim \cite{ButlerT5}: the weighted coframe module $\shW$ on the blown-up manifold, the boundary $E \cong \R \times S^2$, the degenerate Gram $G$ at the boundary with radical spanned by the dying direction $k$, and the one-sided completion that adjoins the plane $(\nO, \nInf)$ and couples it to $k$ through collar data. Throughout, $\Sig$ denotes the singular boundary ($E$ in T5's notation, the tame wall in T4's), and all pointwise statements are at points of $\Sig$. Three imports carry the note; each is stated in the form used, with provenance.

\begin{assembled}[Imported: collar data and the coupling]\label{c:three-routes.imported-collar-data-coupling}
Collar data for $\Sig$ consist of a global spanning section $k$ of the radical line bundle $\shL$ (existing exactly when $w_1(\shL) = 0$, the standing hypothesis), a transversal vector field $u$ with $k(u) = 1$, a bump $\beta$ equal to $1$ on an inner collar, and constants $\mu, \nu \neq 0$. The coupling is the symmetric form $\kappa$ with $\kappa(\alpha, \nO) = \nu\beta\,\alpha(u)$ and $\kappa(\alpha, \nInf) = \mu\beta\,\alpha(u)$ for covectors $\alpha$, vanishing on covector pairs and on the adjoined plane; the ambient form is $h := G \perp H + \kappa$, with $H$ the hyperbolic plane form on $(\nO, \nInf)$, $\nO^2 = \nInf^2 = 0$, $\nO \cdot \nInf = -1$. On the inner collar $\beta = 1$, and every computation of this note is made there. \emph{Provenance:} \cite[Def.~4.1]{ButlerT4}, verbatim up to notation.
\end{assembled}
\begin{proof}
Collar data for $\Sig$ consist of a global spanning section $k$ of the radical line bundle $\shL$ (existing exactly when $w_1(\shL) = 0$, the standing hypothesis), a transversal vector field $u$ with $k(u) = 1$, a bump $\beta$ equal to $1$ on an inner collar, and constants $\mu, \nu \neq 0$. The coupling is the symmetric form $\kappa$ with $\kappa(\alpha, \nO) = \nu\beta\,\alpha(u)$ and $\kappa(\alpha, \nInf) = \mu\beta\,\alpha(u)$ for covectors $\alpha$, vanishing on covector pairs and on the adjoined plane; the ambient form is $h := G \perp H + \kappa$, with $H$ the hyperbolic plane form on $(\nO, \nInf)$, $\nO^2 = \nInf^2 = 0$, $\nO \cdot \nInf = -1$. On the inner collar $\beta = 1$, and every computation of this note is made there. \emph{Provenance:} \cite[Def.~4.1]{ButlerT4}, verbatim up to notation.

\end{proof}

\begin{assembled}[Imported: the boundary Gram and the ladder]\label{c:three-routes.imported-boundary-gram-ladder}
At each point of $\Sig$ the Gram $G$ of $\shW$ is symmetric of corank exactly one, its radical the line $\R k$, with $G(k,k) = 0$ there and $\det G = -r$ off the boundary; for exact Schwarzschild, on the frame $(\rho, dt_{\mathrm{PG}}, \mu_1, \mu_2)$,
\[
G(0) \;=\; \begin{pmatrix} m^2 & -m & 0 & 0\\ -m & 1 & 0 & 0\\ 0 & 0 & -1 & 0\\ 0 & 0 & 0 & -1 \end{pmatrix}, \qquad k = \rho + m\,dt_{\mathrm{PG}}, \qquad m = \sqrt{2M}.
\]
The completed form $h$ is nondegenerate at every point of the physical closure, and the ladder pair $k^{*} = \nO/2\nu$, $kk^{*} + k^{*}k = 1$, $k^{2} = (k^{*})^{2} = 0$, exists exactly over $\Sig$. \emph{Provenance:} \cite[Thms 8.7, 8.12]{ButlerT4}, reproduced with hypotheses as \cite[Facts B.1, B.2]{ButlerT5}.
\end{assembled}
\begin{proof}
At each point of $\Sig$ the Gram $G$ of $\shW$ is symmetric of corank exactly one, its radical the line $\R k$, with $G(k,k) = 0$ there and $\det G = -r$ off the boundary; for exact Schwarzschild, on the frame $(\rho, dt_{\mathrm{PG}}, \mu_1, \mu_2)$,
\[
G(0) \;=\; \begin{pmatrix} m^2 & -m & 0 & 0\\ -m & 1 & 0 & 0\\ 0 & 0 & -1 & 0\\ 0 & 0 & 0 & -1 \end{pmatrix}, \qquad k = \rho + m\,dt_{\mathrm{PG}}, \qquad m = \sqrt{2M}.
\]
The completed form $h$ is nondegenerate at every point of the physical closure, and the ladder pair $k^{*} = \nO/2\nu$, $kk^{*} + k^{*}k = 1$, $k^{2} = (k^{*})^{2} = 0$, exists exactly over $\Sig$. \emph{Provenance:} \cite[Thms 8.7, 8.12]{ButlerT4} (\ref{c:sheaf.schwarzschild-degeneration-boundary-tame}, \ref{c:sheaf.one-sided-completion-singularity}), reproduced with hypotheses as \cite[Facts B.1, B.2]{ButlerT5}.

\end{proof}

\phantomsection\label{c:note.three-routes.setting-setting-what-imported-2}
The three sections that follow are independent of one another. the section ``The metric-level sheaf'' concerns what T5's Proposition 9.6 calls its definitional antecedent; the section ``The boundary no-go'' closes the open item recorded in the scope paragraph of T5's Fact 5.4; the section ``The hopping exclusion'' proves the class-relative form of the no-hopping statement that T5's Remark 7.33(b) requests. Every computational identity below is verified in exact rational arithmetic by the deposited script (the section ``Machine verification'').

\section{The metric-level sheaf}

\phantomsection\label{c:note.three-routes.setting-metric-level-sheaf}
T5's Proposition 9.6 proves that the composite of the pointwise descents is the pinching $\EN$ and that its range is the joint-diagonal algebra; what it leaves as an antecedent is that the metric-level reading of a \emph{region} of the ledger is the simultaneous performance of the pointwise descents. This section shows that the antecedent is the statement that the metric level is \emph{local}: the pointwise metric levels assemble into a genuine sheaf on $\Sig$, that sheaf's sections are the compatible families, its finite truncations are exactly $\operatorname{range}(\EN)$, and the rival reading is disqualified on sheaf-theoretic grounds alone, not by preference. One point deserves emphasis before the definitions: the risk named at the outset, that the pointwise even algebras might fail to glue, resolves positively, and the obstruction that would have prevented gluing is precisely the standing hypothesis the completion already carries.

\begin{definition}[The ladder bundle, and the even subbundle]\label{c:obj.three-routes.ladder-bundle-even-subbundle}
With collar data fixed (Fact \ref{c:three-routes.imported-collar-data-coupling}), the sections $k$ and $k^{*} = \nO/2\nu$ of the completed Clifford sheaf are global over $\Sig$, and at each $x \in \Sig$ they generate the fibre algebra $\Aa_x \cong \Mat_2(\R)$-type mode factor of \cite[Cor.~4.5]{ButlerT4}. Write $\sigma(x) := k^{*}(x)k(x) - k(x)k^{*}(x)$, a global section with $\sigma(x)^2 = 1$, and $\alpha_x := \operatorname{Ad}(\sigma(x))$ for the fibrewise parity. The \emph{ladder bundle} $\Aa \to \Sig$ is the bundle with fibres $\Aa_x$; the \emph{even subbundle} $\Aa^{+} \subseteq \Aa$ has fibres the fixed algebras $\Aa_x^{+} = \{a \in \Aa_x : \alpha_x(a) = a\}$. The \emph{metric-level presheaf} assigns to each open $U \subseteq \Sig$ the algebra
\[
\mathfrak{M}(U) \;:=\; \Gamma(U, \Aa^{+}) \;=\; \bigl\{\, s : U \to \Aa^{+} \ \text{smooth},\ s(x) \in \Aa_x^{+}\ \forall x \in U \,\bigr\},
\]
with the restriction maps of sections.
\end{definition}

\begin{proposition}[It is a sheaf, and the risk resolves into the standing hypothesis]\label{c:three-routes.it-sheaf-risk-resolves-into-standing}
$\mathfrak{M}$ is a sheaf of algebras on $\Sig$. Existence of the untwisted ladder bundle, from which $\mathfrak{M}$ is built, is equivalent to $w_1(\shL) = 0$: the fibrewise condition defining $\Aa^{+}$ uses the parity section $\sigma = k^{*}k - kk^{*}$, $\sigma$ is built from the ladder $(k, k^{*})$, and the ladder exists globally exactly when collar data do, which is $w_1(\shL) = 0$. $\mathfrak{M}$ itself, built from $\operatorname{Ad}\sigma$, which is insensitive to the sign of $k$, glues on a twisted ladder bundle as well and does not see the twist.
\end{proposition}
\begin{proof}
The statement was corrected on 2026-10-23: the untwisted ladder bundle exists exactly when $w_1(\shL) = 0$, but $\mathfrak{M}$, built from $\operatorname{Ad}\sigma$, glues on a twisted ladder bundle as well (as the even-condition clause below already notes), so the closing claim that $\Aa^{+}$ exists only in the $\shL$-twisted sense holds for the ladder bundle and not for its even part. Locality and gluing are inherited from the sheaf of sections of $\Aa$, because membership in $\mathfrak{M}(U)$ is a fibrewise condition: a family of sections $s_i \in \mathfrak{M}(U_i)$ agreeing on overlaps glues to a unique section $s$ of $\Aa$, and $s(x) \in \Aa_x^{+}$ holds at each $x$ because it holds for whichever $s_i$ is defined there. The restriction maps exist because the ladder bundle is a bundle of algebras with smooth ($C^\infty$-linear) structure maps, the machinery's linearity (the note's Fact 1.3(ii), \ref{c:sheaf.function-linearity-tame-constructions}): nothing in the construction differentiates, so restriction is restriction of coefficient functions. The equivalence with $w_1(\shL) = 0$ is the chain stated: without a global $k$ the parity $\sigma$ is defined only up to the sign of $k$, the even condition is still well defined ($\alpha_x$ is insensitive to the sign, $\sigma \mapsto -\sigma$ giving the same automorphism), but the ladder bundle itself, and with it $\Aa^{+}$, exists only in the $\shL$-twisted sense of \cite[\S 4]{ButlerT4}; the untwisted statement used by the chain is exactly the collar-data case (\ref{c:three-routes.imported-collar-data-coupling}, collar data existing exactly when $w_1(\shL) = 0$).

\end{proof}

\begin{proposition}[The stalks are the pointwise metric levels]\label{c:three-routes.stalks-are-pointwise-metric-levels}
For each $x \in \Sig$ the stalk of $\mathfrak{M}$ at $x$ is $C^\infty_x \otimes \Aa_x^{+}$: germs of families valued in the fibre's even part. In particular the stalk's algebra fibre is $\Aa_x^{+}$, the image of the point's own descent, which is the pointwise metric level of \cite[Prop.~5.5, Thm.~5.7]{ButlerT5}.
\end{proposition}
\begin{proof}
$\Aa^{+}$ is a locally trivial subbundle (its fibrewise projection is $\tfrac12(\id + \alpha_x)$, the point's descent (\ref{c:singularity.descent-parity-average}), a smooth idempotent bundle map, again by the machinery's linearity, the note's Fact 1.3(ii), \ref{c:sheaf.function-linearity-tame-constructions}), and the stalk of the sheaf of sections of a bundle is the germs of its local sections.

\end{proof}

\begin{theorem}[The join is disqualified: wrong stalks]\label{c:three-routes.join-disqualified-wrong-stalks}
For $U$ with at least two points define the join $J(U) := \bigl\langle \bigcup_{x \in U} \operatorname{Fix}(\operatorname{Ad}\sigma(x)) \bigr\rangle$, the algebra generated inside the $U$-ledger by the pointwise fixed algebras. Then: (i) $J(U)$ is the whole ledger, so the two-ledger discrepancy computed against it vanishes identically; and (ii) $J$ is the section object of no sheaf on $\Sig$ whose stalks are the pointwise metric levels: any presheaf with $J$ as sections contains, at every point $x$, elements odd at $x$, so its stalk data at $x$ properly exceeds $\Aa_x^{+}$.
\end{theorem}
\begin{proof}
(i) For $y \neq x$ the odd generator $k(y)$ is fixed by $\operatorname{Ad}(\sigma(x))$, since the two points' ladders commute up to the grading and $\sigma(x)$ acts trivially on the $y$-factor's occupation-preserving part; concretely, in the joint representation $\operatorname{Ad}(\sigma_x)(v_y) = v_y$. So $J(U)$ contains every generator $k(y)$, $k^{*}(y)$ for every $y \in U$ (each lies in the fixed algebra of any \emph{other} point), and these generate the ledger. (ii) The element $v_y \in J(U)$ is odd at $y$: $\operatorname{Ad}(\sigma_y)(v_y) = -v_y$. A sheaf on $\Sig$ with stalks $C^\infty_x \otimes \Aa_x^{+}$, the pointwise metric levels of Proposition \ref{c:three-routes.stalks-are-pointwise-metric-levels}, has every section even at every point of its domain, and $v_y$ is not; so no restriction structure can present $J$ as such a sheaf's sections. The script verifies both computations on two modes: $\operatorname{Ad}(\sigma_x)$ fixes $v_y$, $\operatorname{Ad}(\sigma_y)$ negates it.

\end{proof}

\begin{remark}[What this settles, and the exact residue]\label{c:note.three-routes.what-settles-exact-residue}
The two candidate readings of ``the metric level over $U$'' are now separated by a structural property, not a preference: the meet is the section algebra of a genuine sheaf whose stalks are the proved pointwise metric levels; the join has the wrong stalks and is no sheaf at all. What remains assumed is exactly this: \emph{that the metric level over a region is a local notion}, an assignment restricting along inclusions with the pointwise levels as stalks. Grant that, and T5's Proposition 9.6 antecedent follows, $\EN$ is the metric-level reading, and the two-ledger discrepancy of T5's Steps 5 and 6b is the sheaf's verdict rather than a stipulation. Decline it, and one must hold that reading a region at the metric level is not determined by readings near each of its points, against the pointwise character of everything else in the localisation. The residue is a functoriality axiom, and the note states it as one.
\end{remark}

\section{The boundary no-go}

\phantomsection\label{c:note.three-routes.setting-boundary-no-go}
T5's Fact 5.4 records that the no-c-number-repair theorem of \cite[Thm.~4.2]{ButlerT1} is proved where the ambient form is nondegenerate and the dying direction has a hyperbolic partner, and that at the boundary itself, where the partner is exactly what is missing, only the weaker statement was available: no equivariant value had been \emph{exhibited}. This section closes the gap: none \emph{exists}. The candidate objects are fixed first, then the two lemmas, then the two theorems. Everything is at a point of $\Sig$, on the inner collar where $\beta = 1$, with $(W, G, k)$ as in Fact \ref{c:three-routes.imported-boundary-gram-ladder}; ``vector'' means an element of the dual of $W$, and for a vector $v$ and covector $\alpha$ we write $\alpha(v)$ for the pairing.

\begin{definition}[Candidates, and the base unipotents]\label{c:obj.three-routes.candidates-base-unipotents}
A \emph{candidate value} for the dying direction $k$ is a transversal: a vector $u$ with $k(u) = 1$, exactly the datum a classical reciprocal of $k$ would supply and exactly the datum the coupling construction consumes (Fact \ref{c:three-routes.imported-collar-data-coupling}). The candidates form an affine space over $\ker k := \{v : k(v) = 0\}$, of dimension three. For $\lambda \in \ker k$ define the \emph{base unipotent} $g_\lambda : W \to W$, $g_\lambda(\alpha) := \alpha + \alpha(\lambda)\,k$.
\end{definition}

\begin{lemma}[The base unipotents, and their action on candidates]\label{c:three-routes.base-unipotents-their-action-candidates}
Each $g_\lambda$ is a $G$-isometry of $W$ at the boundary, fixes $k$, and $g_\lambda g_{\lambda'} = g_{\lambda + \lambda'}$. The dual action on vectors is $u \mapsto u - k(u)\,\lambda$; on candidates it is the translation $u \mapsto u - \lambda$, a free and transitive action of $(\ker k, +)$.
\end{lemma}
\begin{proof}
$G(g_\lambda\alpha, g_\lambda\alpha') = G(\alpha, \alpha') + \alpha'(\lambda)G(\alpha, k) + \alpha(\lambda)G(k, \alpha') + \alpha(\lambda)\alpha'(\lambda)G(k,k) = G(\alpha, \alpha')$, every correction term vanishing because $k$ spans the radical (\ref{c:three-routes.imported-boundary-gram-ladder}). $g_\lambda(k) = k + k(\lambda)k = k$. The composition law is immediate, $g_\lambda^{-1} = g_{-\lambda}$. For the dual action, $(g_\lambda^{-\mathsf{T}}u)(\alpha) = (g_\lambda^{-1}\alpha)(u) = \alpha(u) - \alpha(\lambda)k(u) = \alpha(u - k(u)\lambda)$; at $k(u) = 1$ this is translation by $-\lambda$, which is free (no fixed vector for $\lambda \neq 0$) and transitive (any two candidates differ by an element of $\ker k$).

\end{proof}

\begin{lemma}[The radical gauge closes the solve]\label{c:three-routes.radical-gauge-closes-solve}
For every $v \in \ker k$ there is a covector $\tilde v \in W$ with $G(\alpha, \tilde v) = \alpha(v)$ for all $\alpha$, unique modulo the radical $\R k$; and the normalisation
\[
G(\tilde v, \tilde v) + 2\,\tilde v(u) \;=\; 0
\]
fixes it uniquely, the left side being affine in the gauge parameter with slope $2\,k(u) = 2 \neq 0$.
\end{lemma}
\begin{proof}
The musical map $\alpha \mapsto G(\alpha, \cdot)$ has kernel $\R k$ (the radical, by \ref{c:three-routes.imported-boundary-gram-ladder}) and image contained in the vectors killed by pairing with $k$, i.e.\ in $\ker k$ regarded inside the bidual; its rank is three, the dimension of $\ker k$, so the image is exactly $\ker k$: solvability, with ambiguity the kernel. Under $\tilde v \mapsto \tilde v + tk$ the quantity $G(\tilde v, \tilde v)$ is unchanged ($k$ radical) and $\tilde v(u)$ shifts by $t\,k(u) = t$, so the displayed normalisation is met at exactly one $t$.

\end{proof}

\begin{theorem}[Change of transversal is a canonical isomorphism of completions]\label{c:three-routes.change-transversal-canonical-isomorphism-completions}
Fix collar constants $\mu, \nu$ and let $h_u$ denote the ambient form built from the transversal $u$ (Fact \ref{c:three-routes.imported-collar-data-coupling}). For $v \in \ker k$ let $\tilde v$ be as in Lemma \ref{c:three-routes.radical-gauge-closes-solve}, and define the unipotent map $\Phi_v$ on $T = W \oplus \operatorname{span}(\nO, \nInf)$:
\[
\Phi_v\big|_W = \id, \qquad \Phi_v(\nO) = \nO + \nu\,\tilde v, \qquad \Phi_v(\nInf) = \nInf + \mu\,\tilde v .
\]
Then $h_u(\Phi_v x, \Phi_v y) = h_{u+v}(x, y)$ for all $x, y$: $\Phi_v$ is an isometry $(T, h_{u+v}) \to (T, h_u)$ restricting to the identity on $W$, fixing $k$, and inducing an isomorphism of the completed Clifford algebras that carries ladder to ladder. The completion is therefore independent of the collar transversal, canonically over the degenerate base.
\end{theorem}
\begin{proof}
Five pairings of the ambient form of \ref{c:three-routes.imported-collar-data-coupling}. On $W \times W$ both sides are $G$. For $\alpha \in W$ against $\nO$: $h_u(\alpha, \nO + \nu\tilde v) = \nu\alpha(u) + \nu\,G(\alpha, \tilde v) = \nu\alpha(u) + \nu\alpha(v) = \nu\alpha(u + v) = h_{u+v}(\alpha, \nO)$, and identically with $\mu$ against $\nInf$. For $\nO$ against itself: $h_u(\nO + \nu\tilde v, \nO + \nu\tilde v) = 2\nu\,\kappa_u(\tilde v, \nO) + \nu^2 G(\tilde v, \tilde v) = 2\nu^2\tilde v(u) + \nu^2 G(\tilde v, \tilde v) = \nu^2\bigl(G(\tilde v,\tilde v) + 2\tilde v(u)\bigr) = 0$ by the normalisation of \ref{c:three-routes.radical-gauge-closes-solve}, matching $h_{u+v}(\nO, \nO) = 0$; identically with $\mu^2$ for $\nInf$. For the cross pairing: $h_u(\Phi\nO, \Phi\nInf) = -1 + \mu\nu\,\tilde v(u) + \nu\mu\,\tilde v(u) + \mu\nu\,G(\tilde v, \tilde v) = -1 + \mu\nu\bigl(G(\tilde v, \tilde v) + 2\tilde v(u)\bigr) = -1 = h_{u+v}(\nO, \nInf)$. So $\Phi_v$ is an isometry of the stated quadratic modules; an isometry of quadratic modules induces an isomorphism of their Clifford algebras (\ref{c:std-algebra.functoriality}), and $\Phi_v$ fixes $k$ and maps the $h_{u+v}$-ladder $(k, \nO/2\nu)$ to a conjugated $h_u$-CAR pair for the same $k$, the anticommutators being preserved. The map is unipotent, $(\Phi_v - \id)$ having image in $W$ and square zero, so it is invertible with unipotent inverse; canonicity is the normalised uniqueness of $\tilde v$. The script verifies the matrix identity $\Phi^{\mathsf{T}} h_u \Phi = h_{u+v}$ exactly on the Schwarzschild boundary data, for two independent choices of $v$, with a wrong-normalisation control that fails.

\end{proof}

\begin{theorem}[No equivariant value at the boundary]\label{c:three-routes.no-equivariant-value-boundary}
For each base unipotent $g_\lambda$ let $\hat g_\lambda := g_\lambda \oplus \id$ on $T$. Then
\[
\theta_\lambda \;:=\; \hat g_\lambda \circ \Phi_\lambda^{-1}
\]
is an automorphism of the fixed completed structure $(T, h_u)$, restricting to $g_\lambda$ on $W$ and fixing $k$; its induced dual action on candidate values is the translation $u' \mapsto u' - \lambda$. Consequently any assignment of a classical partner value at the boundary that is natural under automorphisms of the completed structure would be a fixed point of a free translation action of $(\ker k, +)$ on the candidates, and none exists: \emph{no equivariant classical value repairs the dying direction at $\Sig$ itself}.
\end{theorem}
\begin{proof}
First, $\hat g_\lambda$ is an isometry $(T, h_{u+\lambda}) \to (T, h_u)$: on $W \times W$ it preserves $G$ (Lemma \ref{c:three-routes.base-unipotents-their-action-candidates}); on the adjoined plane it is the identity; and for the coupling of \ref{c:three-routes.imported-collar-data-coupling}, $h_u(\hat g\alpha, \nO) = \nu(g\alpha)(u) = \nu\alpha(u) + \nu\alpha(\lambda)k(u) = \nu\alpha(u + \lambda) = h_{u+\lambda}(\alpha, \nO)$, identically with $\mu$. Composing with $\Phi_\lambda^{-1} : (T, h_u) \to (T, h_{u+\lambda})$ (\ref{c:three-routes.change-transversal-canonical-isomorphism-completions}) gives $\theta_\lambda : (T, h_u) \to (T, h_u)$, an automorphism. Its restriction to $W$ is $g_\lambda$, since $\Phi_\lambda$ is the identity there, and both factors fix $k$. The induced action on vectors is dual to the action on $W$, hence is $g_\lambda$'s: translation by $-\lambda$ on candidates (Lemma \ref{c:three-routes.base-unipotents-their-action-candidates}). A natural value assignment, one candidate per completed structure commuting with automorphisms, must satisfy $\hat u = \hat u - \lambda$ for every $\lambda \in \ker k$, impossible for any $\lambda \neq 0$. The script verifies $\theta_\lambda^{\mathsf{T}} h_u\, \theta_\lambda = h_u$ exactly, that $\theta_\lambda$ is not the identity, that it fixes $k$, and the translation action on two sample candidates.

\end{proof}

\begin{remark}[Scope, and what this discharges]\label{c:note.three-routes.scope-what-discharges}
The theorems are pointwise on $\Sig$ with every formula rational in the collar data, hence smooth in the point, and so hold over every open piece of $\Sig$; the collar constants $\mu, \nu$ and the bump are held fixed, their variation being a rescaling of the adjoined plane recorded in \cite[\S 4]{ButlerT4} and orthogonal to the question. Within that scope, the open item of T5 Fact 5.4's scope paragraph is discharged: the mechanism of \cite[Thm.~4.2]{ButlerT1} is not merely portable but ported, the affine candidate space and the fixed-point-free unipotent action both realised at the degenerate point, with the completion's own transversal-independence (Theorem \ref{c:three-routes.change-transversal-canonical-isomorphism-completions}) supplying the equivariance that the sketch could only name. What remains for T4 is transcription: restating Theorems \ref{c:three-routes.change-transversal-canonical-isomorphism-completions} and \ref{c:three-routes.no-equivariant-value-boundary} as statements about its sheaves, for which every ingredient above is already sheaf-ready. Upon that transcription, T5's abstract sentence on classical repair may drop its scope qualifier.
\end{remark}

\section{The hopping exclusion}

\phantomsection\label{c:note.three-routes.setting-hopping-exclusion}
T5's Premise 7.30, ledger locality, has two halves; the diagonality half was argued there from what the localisation theorems fail to supply, and T5's Remark 7.33(b) records the gap: ``a no-go on inter-point hopping over $E$ would settle it.'' This section proves the class-relative form of that no-go, and is explicit about which class. The result upgrades the premise's status from \emph{unsupplied} to \emph{excluded within the class of supplied generators}; the residue, stated in Remark \ref{c:note.three-routes.true-residual-differential-dynamics}, is a genuine dynamical question and not bookkeeping.

\begin{lemma}[Function-linear maps preserve supports]\label{c:three-routes.function-linear-maps-preserve-supports}
Let $\mathcal{F}$ be a sheaf of $C^\infty_\Sig$-modules and $D : \mathcal{F}(U) \to \mathcal{F}(U)$ a $C^\infty(U)$-linear map. Then $\supp(Ds) \subseteq \supp(s)$ for every section $s$.
\end{lemma}
\begin{proof}
Let $f \in C^\infty(U)$ equal $1$ on a neighbourhood of $\supp(s)$. Then $s = fs$, so $Ds = D(fs) = f\,Ds$ and $\supp(Ds) \subseteq \supp(f)$. Bumps with $\supp(f)$ inside any prescribed closed neighbourhood of $\supp(s)$ exist on $\Sig \cong \R \times S^2$; intersecting over them gives $\supp(Ds) \subseteq \supp(s)$.

\end{proof}

\begin{proposition}[Support-preserving families quantise without hopping]\label{c:three-routes.support-preserving-families-quantise-without-hopping}
Let $a$ be a $C^\infty$-linear endomorphism of the ladder bundle, hence by Lemma \ref{c:three-routes.function-linear-maps-preserve-supports} a bundle map, acting fibrewise as $a_x : \Aa_x \to \Aa_x$. On any finite truncation $F = \{x_1, \dots, x_N\}$ its second quantisation is the sum of single-point bilinears
\[
d\Gamma(a) \;=\; \sum_{j} \psi^{\dagger}_{x_j}\, a_{x_j}\, \psi_{x_j},
\]
which commutes with every occupation number $N_{x_j}$: no such generator moves occupation between points. By contrast a one-particle map with a nonzero cross component $a_{x y}$, $x \neq y$, quantises to a generator containing $\psi^{\dagger}_x \psi_y$, which maps the occupation sector $(N_x, N_y) = (0,1)$ onto $(1,0)$.
\end{proposition}
\begin{proof}
The block-diagonal form of $d\Gamma(a)$ is the definition of second quantisation applied to a bundle map (by \ref{c:three-routes.function-linear-maps-preserve-supports}), each summand supported at one point; each summand is a polynomial in that point's ladder and commutes with every $N_{x_j}$, its own included when $a_{x_j}$ is expressed in the number basis, and the others trivially. The cross-component statement is the elementary computation $\psi^{\dagger}_x\psi_y\,|0_x 1_y\rangle = |1_x 0_y\rangle$; the script performs both computations as exact matrix identities on two modes, with the controlled-$Z$ generator $\sigma_x\sigma_y$ verified to conserve both occupations, since cross-point \emph{diagonal} products are not hopping and were never excluded.

\end{proof}

\begin{theorem}[The class-relative no-go]\label{c:three-routes.class-relative-no-go}
Every generator the tame machinery supplies is $C^\infty$-linear (the note's Fact 1.3(ii), \ref{c:sheaf.function-linearity-tame-constructions}: the machinery is algebra-linear in the coefficient functions and never differentiates them). Hence, by Lemma \ref{c:three-routes.function-linear-maps-preserve-supports} and Proposition \ref{c:three-routes.support-preserving-families-quantise-without-hopping}, no generator supplied by the localisation, the completion, or any construction of the physical closure moves occupation from the ladder at one point of $\Sig$ to the ladder at another. Within the supplied class, the diagonality half of T5's Premise 7.30 is a theorem: the available cross-point structure is exhausted by occupation-conserving products such as $\sigma_x\sigma_y$, which is exactly the class the premise permits and the finite-ledger model uses.
\end{theorem}
\begin{proof}
Immediate from the two preceding results once the class membership is granted, and the class membership is the note's Fact 1.3(ii), the machinery's linearity (\ref{c:sheaf.function-linearity-tame-constructions}), verbatim: a supplied generator is a construction of the machinery, every construction is $C^\infty$-linear, a $C^\infty$-linear endomorphism is a bundle map by Lemma \ref{c:three-routes.function-linear-maps-preserve-supports}, and a bundle map quantises without hopping by Proposition \ref{c:three-routes.support-preserving-families-quantise-without-hopping}.

\end{proof}

\begin{remark}[The true residual: differential dynamics]\label{c:note.three-routes.true-residual-differential-dynamics}
What Theorem \ref{c:three-routes.class-relative-no-go} does not and cannot exclude is a generator outside the $C^\infty$-linear class: a first-order differential operator $D$ on sections over $\Sig$, of Dirac type, has $[D, f] = c(df) \neq 0$ and is precisely not a bundle map; its quantisation would hop, and nothing above forbids nature such a dynamics on $E$, the closure merely constructing none. The question T5's Remark 7.33(b) asks therefore has a sharp remaining form: \emph{is the physical dynamics on the singular boundary $C^\infty$-linear over $E$, or does it carry a differential term?} A differential term would be a new geometric structure on $E$ beyond the completed module, and none of the series' constructions produces one; deciding whether that is a theorem (no such structure is forced) or a physical assumption is the genuine successor question, and it belongs beside the closure-completeness programme rather than inside this note.
\end{remark}

\section{Machine verification}

\phantomsection\label{c:note.three-routes.setting-machine-verification}
Every computational identity above is verified by the deposited script \texttt{verify\_three\_routes.py}, in exact rational arithmetic against the series' shared library \texttt{GA.py}, with negative controls. The coverage: the two-mode computations behind Theorem \ref{c:three-routes.join-disqualified-wrong-stalks} ($\operatorname{Ad}(\sigma_x)$ fixing $v_y$, $\operatorname{Ad}(\sigma_y)$ negating it, the meet equal to the joint diagonal and to the range of the pinching, the join containing every generator); the sheaf axioms on a three-point model of the even subbundle (locality, unique gluing from an overlapping cover, the join's stalk defect exhibited); the full the section ``The boundary no-go'' block on the Schwarzschild boundary data of Fact \ref{c:three-routes.imported-boundary-gram-ladder} at two rational masses (radical membership $Gk = 0$, the solve of Lemma \ref{c:three-routes.radical-gauge-closes-solve} with its normalisation and its uniqueness, the intertwiner identity $\Phi^{\mathsf{T}} h_u \Phi = h_{u+v}$ for two independent $v$, the ladder relations in both completions, unipotence, the automorphism identity $\theta^{\mathsf{T}} h_u \theta = h_u$ with $\theta \neq \id$ and $\theta(k) = k$, the translation action on two candidates, and a wrong-normalisation control that breaks the intertwiner); and the the section ``The hopping exclusion'' block (the finite support lemma as the exact statement that commuting with the point projectors forces block diagonality, the hopping matrix failing it, $d\Gamma$ of a block map conserving both occupations, $d\Gamma$ of a hop moving $(0,1)$ to $(1,0)$, and the controlled-$Z$ generator conserving both). The script exits nonzero on any failure; no floating point is used anywhere.

\section{Descent, parity and states}

\begin{theorem}[On a finite truncation, meet of the big fibres equals join of the small fibres equals $\operatorname{range}(E_F)$]\label{c:singularity.finite-truncation-meet-big-fibres-equals}
For finite $F = \{x_1, \dots, x_N\} \subseteq E$,
\[
\bigcap_{j} B_j \;=\; \Bigl\langle \bigcup_j \Aa_{x_j}^{+} \Bigr\rangle \;=\; \operatorname{range}(E_F),
\]
the joint-number-diagonal algebra, of dimension $2^N$: the meet of the big fibres and the join of the small fibres coincide, and both are the algebra of compatible evaluations of $\mathfrak{M}$ at the points of $F$. Here $\langle\,\cdot\,\rangle$ denotes the generated subalgebra.
\end{theorem}
\begin{proof}
The meet: an element fixed by every $\operatorname{Ad}\sigma(x_j)$ is diagonal in the joint number basis, and conversely (Proposition \ref{c:singularity.p-reduced-metric-level-reading-many}), and $\operatorname{range}(E_F)$ is that algebra, of dimension $2^N$. The join: the $\sigma(x_j)$ commute, are single-site in the Jordan-Wigner representation (the strings cancel in $v^{*}v - vv^{*}$), and generate the $2^N$ independent monomials $\prod_{j \in S}\sigma(x_j)$, $S \subseteq F$, each joint-diagonal; so the join is a subalgebra of the meet of the same dimension. Compatible evaluations of $\mathfrak{M}$ at the points of $F$ are exactly tuples in $\prod_j \Aa_{x_j}^{+}$ (every tuple occurs, by bump functions), and the algebra they generate is the join.

\end{proof}

\begin{remark}[What the antecedent is, stated plainly]\label{c:note.singularity.what-antecedent-stated-plainly}
Theorems \ref{c:singularity.finite-truncation-meet-big-fibres-equals} and \emph{The net assembly is a net, and it is everything} [T5, Thm C.4] quantify the two readings of ``the metric level over a region'' and do not decide between them. Both are determined by readings near the region's points; they differ in how the pointwise levels are assembled. The contravariant assembly (the sheaf, a limit: an element must survive every point's descent) gives $\operatorname{range}(\EN)$, the joint-diagonal algebra, and with it the nonzero two-level discrepancy of Steps 5 and 6b. The covariant assembly (the net, a colimit: the algebra generated by what each point's descent retains of the whole mode algebra) gives the whole mode algebra and a discrepancy of zero; it is the assembly an algebraic quantum field theory would adopt for its algebras of regions. The chain adopts the former, and that adoption is Proposition \ref{c:singularity.p-reduced-metric-level-reading-many}'s antecedent: the sheaf assembly with the small fibres as stalks is the metric level if and only if the metric level over a region is the meet, so the antecedent is restated by this appendix, not reduced. What the appendix adds is the structural content of the choice: on the sheaf side the reading is a genuine sheaf whose finite truncations are $\operatorname{range}(E_F)$, on the net side it is a net whose finite truncations are the full mode algebras, and the discrepancy of Theorem \emph{The deficit theorem: the cost of descent, exactly} [T5, Thm 9.3] is precisely the cost of reading a region as a limit rather than a colimit. The choice is made in Definition \emph{The two fibres, and the two assemblies} [T5, Def C.1], before any sheaf theory runs, and it is a choice about which algebra ``the pointwise metric level'' names: the even part of the mode factor, or the even-in-that-mode part of the boundary mode space.
\end{remark}

\section{Results used from elsewhere}
The proofs above use the following results of the series, stated briefly; each is proved in the paper named.

\begin{fact}[Function-linearity of the tame constructions]\label{c:sheaf.function-linearity-tame-constructions}
All the tame constructions are linear over smooth functions and stay so under composition and the standard operations, so every endomorphism they supply is tensorial.
\emph{Source:} T4, Prop 9.17 (\href{https://doi.org/10.5281/zenodo.22721669}{doi:10.5281/zenodo.22721669}).
\end{fact}

\begin{fact}[One-sided completion at the singularity]\label{c:sheaf.one-sided-completion-singularity}
At the blown-up Schwarzschild singularity, coupling an adjoined hyperbolic plane to the dying direction with $\mu\nu < 0$ and bump inside the horizon gives an ambient form of constant signature $(2,4)$ everywhere, boundary included, with no shadow wall and a ladder pair existing exactly over the boundary.
\emph{Source:} T4, Thm 8.12 (\href{https://doi.org/10.5281/zenodo.22721669}{doi:10.5281/zenodo.22721669}).
\end{fact}

\begin{fact}[The Schwarzschild degeneration is boundary-tame in $\shW$]\label{c:sheaf.schwarzschild-degeneration-boundary-tame}
In the weighted module at the blown-up Schwarzschild singularity, the derived cometric of the Painlev\'e-Gullstrand gauge field extends smoothly to the boundary and degenerates there with corank exactly one and to exactly first order, with radical $k = \rho + m\,\dtPG$ annihilating the infalling congruence.
\emph{Source:} T4, Thm 8.7 (\href{https://doi.org/10.5281/zenodo.22721669}{doi:10.5281/zenodo.22721669}).
\end{fact}

\begin{fact}[The descent is the parity average]\label{c:singularity.descent-parity-average}
The descent is the parity average $E(x) = \tfrac12\bigl(x + \alpha(x)\bigr)$, the unique conditional expectation onto the even subalgebra, with kernel the odd ladder sector.
\emph{Source:} T5, Thm 5.7 (\href{https://doi.org/10.5281/zenodo.22721790}{doi:10.5281/zenodo.22721790}).
\end{fact}

\begin{fact}[(P$_{\mathrm{loc}}$), reduced: the metric-level reading on many modes]\label{c:singularity.p-reduced-metric-level-reading-many}
The mode-wise descents commute and compose to the pinching $\EN$, so if the metric level is read as all pointwise descents at once, it is the joint-number-diagonal algebra.
\emph{Source:} T5, Prop 9.6 (\href{https://doi.org/10.5281/zenodo.22721790}{doi:10.5281/zenodo.22721790}).
\end{fact}

\begin{fact}[Isometries extend to algebra homomorphisms]\label{c:std-algebra.functoriality}
An isometry $f$ extends uniquely to an algebra homomorphism $\Cl(f)$, injective when $f$ is; $\mathrm{O}(V, Q)$ acts by automorphisms.
\emph{Source:} the algebra core of the Transport Lab, proved there with an exact check; standard, see P. Lounesto, Clifford Algebras and Spinors, 2nd ed., London Mathematical Society Lecture Note Series 286, Cambridge University Press (2001), Ch. 14 (definitions of the Clifford algebra).
\end{fact}

\begin{thebibliography}{99}
\bibitem{ButlerT4} L.~Butler,
\emph{Clifford Sheaves over Degenerate Metrics: Nilpotent Ideals, One-Sided Completion, and the Exceptional Algebra},
Transport Papers T4 (2026), doi:\href{https://doi.org/10.5281/zenodo.22721669}{10.5281/zenodo.22721669}.
\bibitem{ButlerT5} L.~Butler,
\emph{Division by Zero at the Gravitational Singularity: The Transport Principle Applied to the Black Hole, and the Information Paradox as a Category Error},
Transport Papers T5 (2026), doi:\href{https://doi.org/10.5281/zenodo.22721790}{10.5281/zenodo.22721790}.
\bibitem{ButlerT1} L.~Butler,
\emph{The Transport Principle: Total Division in Clifford Algebras via Minimal Completion},
Transport Papers T1 (2026).
\end{thebibliography}
\end{document}